\chapter{Nursing Home Care}

\medskip
\noindent\textit{\textbf{Under review.} This chapter is an AI-assisted educational draft; it is not a substitute for professional medical advice.}

\section*{Definition and Overview}
Nursing home care, also called residential aged care, provides accommodation, personal care, nursing, health support, and assistance with daily activities for people who cannot safely live independently. Care should be person-centred, culturally safe, and based on the resident's preferences, abilities, goals, and assessed needs.

\section*{Who May Need It}
People may need residential care because of frailty, dementia, disability, complex medical needs, falls risk, or lack of safe support at home. Assessment should consider home care, respite, palliative care, and supported living as alternatives where appropriate.

\section*{Assessment and Care Planning}
A comprehensive assessment covers mobility, cognition, nutrition, continence, medicines, skin integrity, pain, mood, communication, falls risk, and advance-care preferences. The care plan should be reviewed when the resident's needs change.

\section*{Safety and Quality}
Residents have rights to dignity, privacy, informed consent, safe care, visitors, complaint pathways, and freedom from abuse, neglect, and avoidable restraint. Report sudden deterioration, unexplained injury, medication problems, or suspected abuse promptly.

\section*{Common Health Needs}
Nursing homes commonly manage dementia, diabetes, continence problems, pressure injuries, falls, infections, swallowing difficulty, polypharmacy, and end-of-life symptoms. A GP, nurse, pharmacist, allied-health team, and family or substitute decision-maker may contribute.

\section*{When to Seek Medical Care}
Call 000 for severe breathing difficulty, chest pain, collapse, stroke symptoms, major bleeding, or an immediately life-threatening change. Contact the facility nurse or GP promptly for fever, new confusion, reduced intake, falls, pain, or a sudden change in function.

\section*{Prevention and Self-Care}
Maintain mobility and activity as able, support hydration and nutrition, keep vaccinations current, review medicines regularly, prevent falls and pressure injuries, and document goals of care.

\section*{Expanded clinical context}
Nursing home care is a clinical procedure or health service. Interpretation should account for age, baseline health, medicines, comorbidities, goals, and access to care. This chapter supports discussion with a qualified clinician and does not establish a diagnosis.

\section*{Assessment and decision points}
Clarify the indication, expected benefit, alternatives, consent requirements, preparation, recovery, and possible complications. Review history, allergies, medicines, pregnancy, previous reactions, and support needs. Document the main concern, time course, functional impact, relevant risks, and red flags. Use targeted investigations when their result is likely to change management.

\section*{Management and follow-up}
Safe care includes identity and site checks, appropriate analgesia or anaesthesia, monitoring, written discharge advice, rehabilitation, and planned follow-up. Explain expected improvement, when reassessment is needed, and how follow-up can be accessed. Avoid starting, stopping, or sharing prescription medicines without professional advice.

\section*{Safety-netting}
Seek urgent review for uncontrolled bleeding, severe pain, fever, spreading redness, wound separation, new neurological symptoms, breathing difficulty, or collapse.

\section*{Practical prevention}
Ask who to contact after hours and which medicines or activities should be paused or resumed. Keep written instructions and seek review if the topic recurs, changes, or does not improve as expected.

\section*{Limitations}
This educational expansion should be checked against current local guidance and authoritative topic-specific sources.
\clearpage
